% Project:  TeX PhD Motivation Letter Template
% Author:   Deep Ghuge

% EDIT main.tex directly

\documentclass[11pt]{report}
\usepackage[margin=0.80in]{geometry}
\usepackage{graphicx}

\usepackage{hyperref}

\newcommand{\mediumsize}{\fontsize{15pt}{16pt}\selectfont}

\begin{document}
	\begin{titlepage}
		
		\begin{minipage}[t]{0.90\textwidth}
			\hfill
			\raggedleft
			\textbf{Edwin Varghese} \\
			London, United Kingdom \\
			\today\\
			\href{mailto:edwinshajivarghese@gmail.com}{edwinshajivarghese@gmail.com} \\
			\url{www.linkedin.com/in/edwinvarghese1}
		\end{minipage}
		
		\raggedright \textbf{Dr Marco Fazzi} \\ School of Mathematical and Physical Sciences \\ University of Sheffield \\ Sheffield, UK\\ \href{mailto:m.fazzi@sheffield.ac.uk}{m.fazzi@sheffield.ac.uk}\\
		
		\vspace{0.7em}
		
		\raggedright Dear Dr Marco Fazzi,\\
		\vspace{0.7em}
		
	I am applying for the PhD position exploring Quantum field theories via holography, their moduli spaces, and quantum black holes. I am particularly interested in this project because it presents me with the valuable opportunity to build upon my QFT knowledge from my degree and extend it with advanced methods used in modern high-energy particle theory. My experience as an MSci Physics student, research projects in particle physics theory, neutrino oscillations, plasma modelling, and computational physics more broadly has convinced me that I want a career in this field. I see this PhD as the next logical step for my development as a theoretical particle physicist. \\
		\vspace{0.7em}
		
	During my on-going Master’s project, in which I am deriving the hadronic cross section for $pp \rightarrow W \rightarrow e\nu$, I discovered that I have a keen interest in formal aspects of QFT, and this project has given me new insight into the relationship between theory and phenomenology. Working through electroweak Feynman rules, gamma-matrix algebra, and spinor structures, I developed an appreciation for how the formal structure of quantum field theory and translates into experimentally testable predictions. Implementing phase-space reductions, numerical integration, and Monte Carlo sampling allowed me to consolidate my skills in analytical and computational methods. Through this project, I learned to engage critically with the mathematics, apply computational techniques, and connect theoretical calculations to experimentally meaningful predictions. Building on the theoretical and computational discipline I developed through my major project, I was able to  apply the same analytical rigour to challenges in experimental particle physics.\\
		\vspace{0.7em}
		
		During a six-week research placement at my university, I implemented and trained neural-network-based classifiers to evaluate a novel MCMC convergence metric ($R^*$) and applying the method to simulated datasets from the T2K and NO$\nu$A experiments. I benchmarked $R^*$ against the industry standard Gelman-Rubin ($\hat{R}$) statistic by measuring the metric’s sensitivity to statistical noise, showing that $R^*$ detects smaller deviations between chains and provides complementary global convergence information where  $\hat{R}$ is insensitive. This experience taught me to approach problems systematically, test assumptions, validate methods, and interpret results in detail, skills which I look forward to bringing into my doctoral research.\\
		\vspace{0.7em}
		
		Complementing this, my UKAEA plasma-simulation internship focused on integrating and benchmarking simulation codes for MAST-Upgrade plasmas. I automated the generation of RABBIT input files --- a novel, fast, neutral-beam injection code for plasma simulation --- containing data about the initial state of the plasma, and the plotting of post-simulation outputs with Python. RABBIT was benchmarked against the higher-fidelity NUBEAM/TRANSP modelling chain using the OMFIT workflow manager. Through this internship I learnt how to organise and check large sets of simulation outputs in a systematic way. I also gained experience writing Python scripts with good documentation to ensure the work could be reproduced and shared easily. Having built confidence working with complex simulation tasks, I then undertook projects directly relevant to collider phenomenology at Royal Holloway.\\
		\vspace{0.7em}
		
		My third-year project explored how variations in the Higgs boson mass affect decay channel contributions; I solved formulae numerically using Python and produced logarithmic plots of branching ratios and widths, presenting my results to faculty and peers. I also completed a research review of Higgs searches at the LHC, in which I discussed the relevant production and decay channels, and the impact of higher-order QCD and electroweak corrections on discovery sensitivity. Although these projects are grounded in Standard Model phenomenology, they prepare me well for engaging with more formal theoretical frameworks.\\
		\vspace{0.7em}
		
		Beyond technical skills, I have leadership and communication experience as Academic Officer of the RHUL Physics Society, where I ran LaTeX and Python workshops, organised a weekly PhD seminar series, and coordinated student-led conferences. I enjoy teaching, presenting, and collaborating --- qualities that I believe will serve me well in a research environment where I look forward to exchanging ideas and gaining insights from others.\\
		\vspace{0.7em} 
		
		My experience across machine learning, simulation work and theory-focused projects has convinced me that I want to build my career in theoretical particle physics. My responsibilities in collaborative research environments have taught me the importance of clear documentation, careful validation and open communication. I am eager to develop and consolidate my skills in a new setting where I can explore quantum field theory at a deeper and more formal level, and it will be a pleasure to exchange ideas with and gain insights from researchers working in this field. What draws me most to this PhD is the opportunity to build directly on the quantum field theory I have studied while expanding my understanding of advanced theoretical tools and approaches. I am particularly excited about the prospect of contributing to research that may provide an insight into the fundamental structures of quantum field theories, and I look forward to growing into an independent researcher within this community.\\
		\vspace{0.7em}
		
		
		Thank you for considering my application.\\
		\vspace{5em}
		
		\raggedright Yours sincerely,\\
		\vspace{0.7em}
		\textbf{Edwin Varghese}
		
	\end{titlepage}
\end{document}